\section{Background: Attacks on IoT Devices}\label{sec:background}

This section sets out the starting point for readers new to IoT security: what an attack is, why anyone
attacks a thermostat or an infusion pump, what they gain, and which attacks matter for a detector that
learns on the device.

\subsection{What is an attack?}
An \emph{attack} is a deliberate action that makes a system behave against the wishes of its owner. Security
is usually stated as three properties, and every attack breaks at least one of them:
\begin{itemize}
\item \textbf{Confidentiality}: only authorised parties can read data (broken by data theft or eavesdropping).
\item \textbf{Integrity}: data and behaviour change only in authorised ways (broken by tampering, forged
  messages, or poisoning a model).
\item \textbf{Availability}: the system works when needed (broken by denial of service, or by a detector
  that drowns its operators in false alarms).
\end{itemize}
An attack needs a \emph{vulnerability} (a weakness such as a default password or a firmware bug) and an
\emph{exploit} (the method that uses it). The attacker's \emph{capabilities} (what they can see and
control) and \emph{goal} together form a \emph{threat model}; ours is stated in Section~\ref{sec:threat}.

\subsection{Why IoT devices are attacked}
IoT devices are attractive targets for five reasons. They are \emph{numerous}, so even a rare weakness
yields many victims. They are \emph{weakly protected}: default credentials, rarely patched firmware and
little memory for security software. They are \emph{always on and always connected}, which makes them
useful as relays. They \emph{sit inside trusted networks}, so a compromised camera is a foothold next to
laptops and servers. And some of them \emph{act on the physical world}: locks, pumps, beds and industrial
controllers. The Mirai botnet showed the scale: by scanning for default passwords it recruited a peak of
about 600,000 devices and used them for some of the largest denial-of-service attacks recorded at the
time~\cite{mirai}.

\subsection{What an attacker gains}
Table~\ref{tab:gains} lists the main motives. Most of them need the attacker to stay on the device for a
long time without being noticed, which is exactly why poisoning the device's own detector is worth the
effort.

\begin{table}[t]
\caption{What an attacker can gain from a compromised IoT device}\label{tab:gains}
\centering\footnotesize
\begin{tabular}{@{}p{2.1cm}p{5.9cm}@{}}
\toprule
Goal & Example and gain \\
\midrule
Botnet recruitment & Device floods a target on command; denial of service sold for hire \\
Data theft & Camera footage, health readings, patient identifiers; sold or used for fraud \\
Surveillance & Live audio or video from inside a home, office or ward \\
Extortion & Device or network locked until a ransom is paid \\
Lateral movement & Device used as a stepping stone into the hospital or company network \\
Sabotage & Wrong dose, unlocked door, stopped production line \\
Resource theft & Device used for cryptocurrency mining or as an anonymising proxy \\
\bottomrule
\end{tabular}
\end{table}

\subsection{Types of attack}
Table~\ref{tab:attacks} groups the attacks relevant to this paper. The first group attacks the device or
the network; it is what an intrusion detector is built to catch. Our evaluation uses real traffic of this
group from N-BaIoT~\cite{nbaiot}: the Mirai and BASHLITE botnets scanning for victims and flooding
targets with ACK, SYN, UDP, TCP and junk traffic. The second group attacks the \emph{detector itself}; it
is the subject of this paper.

\begin{table}[t]
\caption{Attack types. The second group targets the learning detector}\label{tab:attacks}
\centering\footnotesize
\begin{tabular}{@{}p{2.1cm}p{5.9cm}@{}}
\toprule
Attack & What happens \\
\midrule
\multicolumn{2}{@{}l}{\emph{On the device and network}} \\
Scanning & Probing addresses and ports for vulnerable devices (Mirai, BASHLITE scan) \\
Credential guessing & Trying default or common passwords \\
Exploitation & Using a firmware bug to run the attacker's code \\
Command and control & Infected device contacts the attacker for instructions \\
Flooding (DDoS) & ACK, SYN, UDP, TCP or junk floods at a target \\
Exfiltration & Quietly sending stolen data out \\
Spoofing, replay & Forging or re-sending messages as another device \\
\midrule
\multicolumn{2}{@{}l}{\emph{On the learning detector}} \\
Evasion & Shaping attack traffic to look normal to a fixed detector \\
Poisoning & Feeding data so that the detector \emph{learns} the attack as normal \\
Slow poisoning & Poisoning in steps too small to alarm (the boiling frog) \\
Freeze & Tripping a stop-learning rule so honest drift is never learned \\
Forged peer updates & Sending false updates in the name of a neighbour \\
Byzantine peers & Real neighbours with valid keys sending biased updates \\
\bottomrule
\end{tabular}
\end{table}

\subsection{Detecting intrusions on the device}
A \emph{signature} detector matches traffic against known attacks; it cannot see new ones and needs
constant updates. An \emph{anomaly} detector learns what normal traffic looks like and flags what differs,
so it can catch attacks it has never seen and needs no labelled attack data. That suits IoT, where every
device model has its own normal behaviour. We use the online centroid detector studied by Kloft and
Laskov~\cite{kloft}: a sample is scored by its distance from a learned centre, in units of a learned
spread, and raises an alarm above a calibrated threshold (Section~\ref{sec:model}).

\subsection{Why a detector must keep learning}
Normal traffic drifts. A new gateway, a VPN tunnel, a changed polling rate or a firmware update shifts
every packet statistic. A detector fixed at deployment then flags normal traffic: in our strongest
drift, 78\% of normal traffic alarms on a detector that never learns. So the detector must learn
after deployment, and on a device there is no human to label the data. It must label its own traffic:
whatever does not raise an alarm is learned as normal.

\subsection{Poisoning, slow poisoning and the freeze attack}
Self-labelled learning is exactly what a poisoner exploits. The \emph{slow poisoner} sends traffic just
inside the alarm boundary. It does not alarm, so it is learned; the boundary moves toward it; the next
step can sit a little further out. Repeated, this walks the detector until the real attack is inside
the normal region, the boiling-frog attack studied since ANTIDOTE~\cite{antidote} and Kloft and
Laskov~\cite{kloft}. The natural defence is to stop learning when traffic looks suspicious. That creates
the \emph{freeze attack}: the attacker sends a burst of loud traffic, the defence stops learning, the
attacker leaves, and honest drift continues unlearned, so false alarms climb and stay high. The core
difficulty is that from one device's own stream, slow honest drift and slow poisoning look the same
(Proposition~\ref{prop:frontier} and~\cite{warden}). The rest of this paper is about breaking that tie.
